
We set out to filter Python SFT training data using doctest execution. We found that only
1 in 1,000 Python files in a curated corpus contains an independently-executable doctest ---
a $31\times$ gap below the \texttt{>>>} pattern prevalence. Python library code is designed for installed
environments, not clean subprocess execution, and import isolation rejects the large majority
of AST-parseable doctests at the import stage.

We contribute the first systematic measurement of this gap at corpus scale, validated with
28/28 unit tests across a three-phase filtering pipeline (pattern detection, AST extraction,
subprocess execution) that processes 10,000 files in 129.8 seconds. We establish compile-only
filtering as the practical execution-quality gate for raw Python SFT corpus construction and
design the SFT quality comparison experiment as immediate future work.

The \texttt{>>>} prompt in Python code is documentation intent, not executable reality in clean
subprocess environments --- measuring the gap between the two is the first step toward
principled execution-quality data curation for code LLMs.

